\documentclass[11pt,a4paper]{article}
\usepackage[T1]{fontenc}
\usepackage[utf8]{inputenc}
\usepackage[english]{babel}
\usepackage{amsmath,amssymb}
\usepackage[margin=2cm]{geometry}
\setlength{\parindent}{0pt}
\setlength{\emergencystretch}{3em}
\allowdisplaybreaks

\title{About \texttt{irrational\_equations\_to\_solve.tex}}
\author{}
\date{}
\begin{document}
\maketitle
This note explains what the test data in \texttt{irrational\_equations\_to\_solve.tex} is used for. That file is input for the engine in \texttt{main.c}; it is read by the test \texttt{test\_irrational\_equations.c} in this folder.

\begin{enumerate}
\item \textbf{Purpose.} The file holds 10 requests of the form $\{x \in \mathbb{R} \mid \ldots\}$: the same equations as \texttt{book\_irrational\_equations.tex}, with nothing claimed. The engine must find the solutions.

\item \textbf{What it exercises.} Reading the solutions off the sign table when the comparisons contain square roots of the unknown, and re-reading what was written.

\item \textbf{Variety.} Textbook exercises 755 to 764. The expected solutions, written in the test, are the results printed in the textbook; the engine must write them the same way.

\item \textbf{Expected result.} The line \texttt{PASS: 10 of 10 solution sets found and equal to the expected ones}.

\item \textbf{How to run it.} From the top folder of the repository:
\begin{verbatim}
cc -std=c11 -O2 -Wall -Wextra -o /tmp/test_irrational_equations \
    tests/equations_inequalities/test_irrational_equations.c \
    && /tmp/test_irrational_equations
\end{verbatim}
\end{enumerate}
\end{document}
